\scanpage{14}
\mspage{29}
\deleted{
images dans $S$ sont des parties $S_i$ à la fois ouvertes et fermées, et
elles recouvrent $S$. Si $S$ est connexe ou quasi-compact, il existe donc
un ensemble fini d'indices \deleted{\unsure{tel}} $i$ tels que les $S_i$ recouvrent $S$, soit
$S'$ la réunion des $P_i$ correspondants, \deleted{\unsure{qui est donc tel que}}\inserted{Alors}
$S'\unsure{-|\longrightarrow}S$ est fini étale surjectif, et posant
$P'=P\times_S S'=\underline{\operatorname{Isom}}_{S'\text{-gr}}(H',G')$, on voit que $P'$
a une section sur $S'$, donc il existe un isomorphisme de $S'$-groupes
$H'=H\times_S S'\xrightarrow{\sim}G'=G\times_S S'=D_{S'}(M)$, ce qui prouve que $S'$ splitte $H$. Le
cas où les composantes \deleted{irréductibles}\inserted{connexes} de $S$ sont \deleted{\unsure{connexes}} ouvertes se ramène
aussitôt au cas où $S$ est connexe. -- Remarquons d'ailleurs qu'on peut,
lorsque $H$ est de type fini sur $S$, \deleted{\unsure{connexe}} et de type constant, donner
\unsure{he|le} critère d'isotrivialité suivant (où il n'est plus nécessaire de faire
aucune restriction sur $S$): $R$ est réunion d'une \uline{suite} de parties à la
fois ouvertes et fermées finies sur $S$.

\uline{Lemme} 5.1\unsure{3|2}. Soient $S$ un préschéma localement noethérien et connexe,
$\unsure{P|R}$ un $S$-préschéma constant tordu quasi-isotrivial, $Z$ une partie à la fois
ouverte et fermée de $\unsure{P|R}$, telle qu'il existe un $s\in S$ tel que $Z_s$ soit fini.
Alors $Z$ est fini sur $S$.
\margin{$Z_s$}
\margin{$Z$}

\deleted{Oublions un instant que $S$ est supposé connexe, et}\inserted{Ne supposons pas d'abord $S$ connexe;}
soit $U$ l'ensemble
des $s\in S$ tels que $Z_s$ soit fini, nous allons prouver que $U$ est à la
fois ouvert et fermé, et que $Z|U$ est fini sur $U$. C'est là évidemment
un \deleted{énoncé} énoncé essentiellement équivalent à 5.1\unsure{3|2}. mais qui a l'avantage
d'être ``de nature locale'' sur $S$ pour la topologie étale (disons), ce
qui nous ramène au cas où $\unsure{P|R}$ est constant, i.e. de la forme $\unsure{I|K}_S$, où $I$ est
un ensemble. (NB l'hypothèse localement noethérienne n'est pas perdue
par un changement de base étale; c'est là où on utilise la quasi-isotrivialité
de $P$ sur $S$). On peut de plus supposer $S$ connexe, \deleted{mais alors}
}
\mspage{13}

